\section{引言}

电力、石化、冶金、交通、新能源等工业领域的大型基础设施通常具有投资规模大、服役周期长、设备类型多和运行环境复杂等特点，其设备与结构在长期运行过程中不可避免地出现磨损、腐蚀、松动、裂纹、变形、局部过热以及其他异常状态。及时获取设备状态并发现早期缺陷，是保障重大基础设施安全运行、降低非计划停机风险和实现预测性维护的重要基础。传统人工巡检依赖巡检人员携带检测设备进入现场，通过目视观察、红外测温以及仪器测量等方式完成设备状态检查，在结构规则、作业空间宽敞的条件下仍具有较强的灵活性。然而，随着工业装置规模扩大和安全生产要求提高，人工巡检面临的局限性越来越突出：一方面，高空、狭窄、高温、有毒有害和强电磁区域会直接增加人员暴露风险；另一方面，大范围设备点位多、重复性工作量大，人工巡检难以保持稳定的频率和一致的评价标准。此外，依赖经验的判断方式还容易造成异常发现时机、记录方式和处置优先级的差异。因此，利用无人化装备替代人员进入高风险和难抵达区域，并通过智能感知与数据分析实现连续化、标准化巡检，已经成为工业智能运维的重要发展方向。

无人机（Unmanned Aerial Vehicle, UAV）具有机动灵活、部署速度快、非接触作业以及可搭载多类传感器等特点，可以在不改变被检设备运行状态的条件下获得设备表面、空间几何和热状态信息。已有研究已将无人机应用于输电线路、桥梁、风电、光伏、工业装置以及地下和受限空间等多类场景
\cite{shakhatreh2019survey,floreano2015science,kumar2012opportunities,cai2014survey,kendoul2012survey}。在电力巡检领域，无人机视觉系统能够对杆塔、绝缘子和金具等部件开展图像采集与缺陷识别；在桥梁检测领域，无人机可以以较低成本获取高分辨率影像并辅助裂缝、剥落、锈蚀等病害评估；在光伏和风电场景中，无人机还可结合红外热像实现热斑和温度异常检测
\cite{nguyen2018automatic,miao2020uav,zhouwq2024review,morgenthal2014quality,ham2016visual,dorafshan2018bridge,spencer2019advances,jadin2012recent,usamentiaga2014infrared}。随着自主飞行、计算机视觉和人工智能技术不断发展，无人机巡检已经由“人工驾驶+图像采集”逐步向“自主感知+自动识别+任务决策”演进。

需要指出的是，复杂工业场景并不只是将普通无人机巡检简单地搬到工业厂区。工业现场通常同时具有空间结构复杂、环境状态变化剧烈和巡检对象高度异构等特征。厂房、管廊、储罐、炉体、塔架以及设备连接部件之间存在大量狭窄通道、重复钢结构和遮挡区域，GNSS信号容易受到遮挡、多径效应以及局部干扰影响；高温辐射、强反光、粉尘、烟雾、蒸汽、夜间低照度和雨雾等因素又会引起视觉或红外观测退化；输电设施、冶金装置等场景还可能存在较强的电磁干扰，对无人机通信和部分传感器工作状态产生影响
\cite{scaramuzza2014vision,stoecker2017review,lvyang2019sense,zeng2016wireless,mozaffari2019tutorial}。与此同时，工业设备异常具有明显的长尾性质，正常状态样本远多于异常样本，真实缺陷又具有突发性和低频性，导致传统依赖大量标注数据的监督学习方法难以直接覆盖全部异常类型
\cite{taox2021survey,pang2021deep}。在任务层面，巡检对象通常数量较多且重要程度不同，单架无人机还受到续航时间、载荷能力、机载算力和通信距离等约束，因此巡检任务不能仅以最短飞行路径为目标，而需要综合考虑覆盖范围、观测质量、异常风险、任务时效和能源约束
\cite{zeng2016wireless,zhang2019energy,zhou2018computation}。

由此可见，真正意义上的复杂工业场景无人机自主巡检，并不是按照预先设定的航线完成数据采集，而是一个由环境理解、设备认知和任务决策共同组成的闭环过程。无人机首先需要知道自身所处位置、周围环境以及可安全通行的空间；随后需要从可见光、红外、激光等多模态数据中识别设备、部件和异常状态；最后还应根据巡检结果实时决定是否重新观测、是否改变任务顺序以及是否由其他无人机接替任务。换言之，自主巡检的核心能力可以概括为三个层次：在感知层回答“无人机在哪里、周围有什么”，在认知层回答“设备是否异常、异常是什么”，在决策层回答“下一步应该如何巡检以及多架无人机如何协同”。这一认识也决定了本文的综述主线。

围绕上述主线，本文从三个方面系统梳理复杂工业场景无人机自主巡检的研究现状。首先，从自主感知和信息融合出发，总结可见光、红外、激光雷达、惯性和卫星导航等多源信息的特点，并重点讨论视觉SLAM、视觉--惯性、激光--惯性及多模态融合技术，分析复杂工业环境下的感知退化问题。其次，从设备异常认知出发，梳理传统监督式视觉检测、小目标缺陷检测、缺陷分割以及小样本、无监督、自监督异常检测的发展，并进一步讨论视觉基础模型和多模态异常认知为开放世界工业巡检带来的新机会。最后，从自主决策出发，分析覆盖路径规划、动态重规划、主动感知、任务分配和多无人机协同的发展趋势，并进一步讨论通信、能源和算力约束。通过对不同技术路线进行比较，本文试图从“技术是否先进”进一步转向“是否适用于真实复杂工业现场”的视角，总结当前自主巡检在鲁棒性感知、未知异常认知、状态驱动规划、群体协同以及系统闭环等方面的关键不足。



从研究范围上看，复杂工业场景无人机自主巡检与一般无人机应用、普通机器人导航以及工业视觉检测之间存在交叉，但又不能简单等同。它既需要借鉴机器人领域的SLAM、运动规划和多智能体协同理论，也需要吸收工业视觉中的缺陷检测、异常检测和状态评估方法，还必须考虑无人机平台带来的飞行安全、能源和通信约束。因此，本文不试图对所有无人机技术进行百科式罗列，而是围绕“感知是否可靠、认知是否可信、决策是否闭环”三个问题组织文献。对于已有成熟的飞控和基础通信技术，本文仅作必要介绍；对于直接决定复杂工业自主巡检能力的感知、异常认知和自主决策，则重点梳理技术演进和工程适应性。这样的范围界定能够避免将综述变成传统无人机技术总览，也能够突出复杂工业场景下的特殊科学问题。

\section{复杂工业场景下无人机自主感知与信息融合}

\subsection{复杂工业场景的感知需求与自主定位问题}

复杂工业场景中的感知任务具有明显的多层次特征。对于大范围厂区，首先需要快速建立环境的整体空间关系，掌握设备、道路、管廊和作业区域的相对位置；进入设备附近后，又需要获得足够精细的局部几何和图像信息，以支撑部件定位和缺陷判断；在重复巡检任务中，还希望将本次观测与历史地图、历史图像和设备台账关联起来，从而判断设备状态的变化趋势。因此，自主巡检感知并不等价于单纯的目标检测或单纯的定位，而需要同时满足导航安全感知、巡检对象感知和场景语义理解三个层面的信息需求
\cite{kanellakis2017survey,nex2014uav,yeum2015vision,tao2019digital}。

从导航安全角度看，无人机需要识别建筑物、管道、塔架、线缆以及人员车辆等静态和动态障碍物，并估计自身与这些目标之间的空间关系。与普通室外低空环境相比，工业厂区往往具有高密度的人工结构，大量设备外形规则且存在重复纹理，局部区域还可能存在明显的遮挡。对于视觉系统而言，重复钢结构容易产生错误匹配，弱纹理表面不利于特征跟踪，而强反光又会改变局部亮度分布。对于激光雷达而言，狭窄空间、玻璃或高温辐射环境可能影响有效回波，并在设备边缘产生稀疏或不稳定点云。因此，工业环境对感知系统的要求不是简单追求某一时刻的高精度，而是要求在长期飞行过程中维持稳定的状态估计和可用的地图质量。

从巡检对象角度看，工业设备通常呈现“大结构+小部件+细小缺陷”的层次化特点。例如塔架、储罐和厂房属于大尺度对象，阀门、法兰、螺栓和连接件属于中尺度部件，而裂纹、腐蚀点、局部变色等又属于细粒度异常。无人机在不同距离和视角下获取的数据具有明显的尺度变化，因此巡检感知必须同时处理大尺度结构定位和小目标精细观测。对于电力设施，线路和杆塔的空间几何关系决定了无人机的安全接近路径，而绝缘子、金具和连接件的状态又需要高分辨率图像才能判断
\cite{deng2014unmanned,miao2020uav,vieira2023insplad}. 对于桥梁与大型设备，则需要将局部缺陷映射到整体三维结构中，才能真正回答“缺陷发生在什么位置、与哪个设备部件相关”这一工程问题
\cite{liu2020image,li2022bridge,jiang2022vision}。

从感知手段看，可见光相机仍然是工业无人机巡检最基础的传感器。其优点是成本低、分辨率高、数据形式直观，同时能够提供大量可用于视觉识别和定位的纹理信息。对于裂纹、锈蚀、剥落、异物等表观缺陷，可见光图像通常是最直接的数据源。然而，可见光成像对照度、阴影、曝光以及视角较为敏感，在夜间、背光和强反光条件下容易产生明显退化。红外热像仪则能够提供与温度相关的信息，适用于电气接头、轴承、光伏组件和其他热状态敏感设备的检测
\cite{jadin2012recent,usamentiaga2014infrared}。但红外图像的空间分辨率、温度标定以及环境温度变化都会影响异常判定，其信息也不能完全替代可见光对结构表面细节的描述。激光雷达能够直接提供三维空间点云，在弱光环境和部分无纹理场景中具有优势，可用于结构测量、障碍物感知和自主建图
\cite{zhang2014loam,wujq2021review}；但高质量点云处理需要一定的机载算力，并且传感器成本与安装条件通常高于普通相机。IMU具有高频率和短时间稳定的特点，但存在积分漂移，因此通常需要与视觉、激光或GNSS信息联合使用。

GNSS仍然是大范围室外无人机作业中最重要的绝对定位来源之一。RTK等高精度定位技术可以减少常规卫星定位的绝对位置误差，因此在开阔环境中的预设航线巡检具有较高实用价值。然而，复杂工业现场的大量钢结构、建筑物和高大设备容易造成卫星信号遮挡和多径效应，局部区域还可能存在电磁干扰，使GNSS定位精度明显下降。对于厂房内部、地下空间、隧道和压力管道等环境，GNSS甚至无法提供有效定位信息，必须依赖机载传感器完成相对定位
\cite{scaramuzza2014vision,ozaslan2017autonomous,petrlik2020robust,nil2012review}。因此，GNSS并不是自主巡检唯一的定位方案，而应被视为多源状态估计中的一个信息来源。当GNSS可用时提供全局约束，当GNSS退化时则通过视觉、激光和惯性信息维持局部连续定位。

从技术发展脉络来看，无人机自主定位已经形成由单一传感器视觉里程计向多传感器融合SLAM演进的路线。视觉SLAM通过图像特征、直接像素误差或者二者的组合估计相机运动；激光SLAM利用点云几何关系直接建立运动约束；视觉--惯性和激光--惯性方法则利用不同传感器的时间特性和观测互补性提高系统鲁棒性
\cite{durrant2006simultaneous,cadena2016past,gupta2022slam}. ORB-SLAM系列代表了基于特征的视觉SLAM发展过程，从单目场景扩展到双目、RGB-D和视觉惯性系统
\cite{mur2015orb,mur2017orb,campos2021orb}；LSD-SLAM和SVO等方法分别代表直接法和半直接法，在高帧率视觉里程计方面具有代表性；VINS-Mono则通过紧耦合地利用图像与IMU信息提高了状态估计能力
\cite{qin2018vins,faessler2016autonomous,loianno2017estimation}。这些方法为无人机在GPS拒止环境中自主飞行奠定了基础，但其在复杂工业环境中的适应程度仍与传感器配置和环境状态密切相关。

复杂工业场景的另一个重要特点是环境不是完全静态的。工业厂区存在人员、车辆、吊装设备和移动机械，生产过程还可能带来烟雾、蒸汽和局部热羽流。如果SLAM算法默认整个环境静态，那么动态物体上的特征或点云可能被错误地用于状态估计，进而产生定位漂移。因此，工业自主巡检中的感知系统需要从“静态环境下的准确定位”进一步发展为“动态与退化环境中的鲁棒定位”。这一问题不仅涉及前端特征跟踪和点云匹配，还涉及异常观测检测、动态目标剔除、传感器可信度评估以及多源观测权重调整。

从已有工程验证来看，受限空间自主飞行已经取得了一系列具有代表性的成果。Ozaslan等将自主导航、建图和检测流程用于水电压力钢管与隧道检查，证明视觉自主飞行可以服务于具有狭窄结构和GNSS拒止特点的基础设施巡检
\cite{ozaslan2017autonomous}。Petrlík等在受限地下环境中研究了具有鲁棒感知与自主导航能力的无人机系统，并在复杂环境中验证了系统的可行性
\cite{petrlik2020robust}。Mohta等则在GPS拒止和杂乱环境中验证了无人机快速自主飞行的能力
\cite{mohta2018fast}。这些研究说明基础的自主感知与定位技术已经达到可飞行、可建图和可避障的阶段，但工业现场真正关注的是长期稳定性和任务可靠性，即系统能否在多个传感器同时退化、场景持续变化的情况下保持可用，而不是仅在短时间实验中获得一次成功飞行。

综上，复杂工业场景的感知需求可以进一步归纳为“全局定位、局部建图、目标观测和状态反馈”的组合任务。传统以单一传感器为核心的定位策略已经难以应对工业现场的多因素退化，多源感知和信息融合成为自然的发展方向；与此同时，自主定位与巡检目标感知不应完全割裂，地图中的设备语义、目标异常以及历史状态应该反过来为定位与任务规划提供约束。也就是说，复杂工业自主巡检需要的并非孤立的SLAM模块，而是一个能够感知自身状态、理解环境结构并识别观测可靠性的统一感知体系。



进一步从工业巡检流程看，自主感知不仅承担“发现障碍物”的功能，还承担巡检任务中不同阶段的状态切换。例如，在进入厂区时无人机主要关注全局定位和大范围障碍物分布；接近设备后，系统需要提高局部几何建模精度，并控制无人机与设备之间保持合理的安全距离；当相机进入近距离检测状态后，感知重点又由导航转向目标跟踪和观测质量评估。因此，感知系统实际上同时具有“导航感知”和“任务感知”两类需求。二者如果完全采用相同的传感器和分辨率，会造成不必要的计算和能源消耗，而根据任务阶段动态调整感知策略，则有可能显著提高系统效率。

工业现场的空间几何通常还具有显著的层级性。大型厂房和装置构成全局尺度，中型管廊、塔架和储罐构成区域尺度，阀门、法兰、螺栓和焊缝则处于局部尺度。对于无人机自主定位而言，全局尺度信息有助于维持长时间飞行的地图一致性，局部尺度几何则决定无人机能否安全靠近目标。因此，多尺度地图表示以及全局--局部定位的协同对于复杂工业巡检具有重要意义。现有SLAM研究更多从算法层面优化位姿估计，而工业自主巡检还需要把地图划分为“可飞行空间、设备空间和重点观测空间”，从而为后续路径规划和目标观测直接提供语义约束。

此外，工业现场经常存在长期重复巡检任务。同一设备可能每天、每周或每月被重复观测，如果每次巡检都完全重新构建地图，则会造成计算和数据存储上的重复。更有价值的方式是建立可复用的先验地图，并将当前飞行结果与历史地图进行配准，从而快速判断设备位置变化和环境变化。这也意味着自主感知需要从一次性SLAM逐步向长期定位与长期建图发展。对于大尺度工业设施，地图不仅是导航工具，同时也是设备状态记录和变化检测的空间载体。

\subsection{多模态感知融合与SLAM研究}

多传感器融合的基本思想是利用不同传感器的互补观测提高对系统状态或环境状态的估计能力。按照信息处理层次，常见方法可以划分为数据级、特征级和决策级融合
\cite{panq2003basic,wangj2004survey,hey2000multi}。数据级融合尽可能保留原始观测信息，在理论上能够获得较充分的互补信息，但要求传感器具有较高的一致性，包括时间同步、坐标系统一和外参标定等；特征级融合首先从各传感器提取中间表征，再进行联合估计，在信息量、计算量和鲁棒性之间取得平衡；决策级融合则分别处理不同传感器的信息，最后通过规则、概率模型或学习模型实现结果互证，因此对异构传感器的独立性要求更低。对于工业无人机而言，这三种层次并不是相互排斥的，实际系统往往在状态估计中采用紧耦合融合，在设备识别和异常判断中再采用决策级或特征级多模态融合。

视觉SLAM的研究重点经历了由特征法、直接法到深度特征和语义增强的演变过程。SIFT、SURF和ORB等局部特征解决了图像关键点的提取和匹配问题，RANSAC和多视图几何为运动估计提供了鲁棒基础
\cite{lowe2004distinctive,rublee2011orb,fischler1981random,hartley2003multiple}。ORB-SLAM、ORB-SLAM2和ORB-SLAM3在这一框架下逐渐增强对不同相机形式、视觉惯性以及多地图场景的支持
\cite{mur2015orb,mur2017orb,campos2021orb}。与特征法相比，直接法不依赖显式的关键点匹配，而是直接最小化像素重投影误差，在纹理较弱但局部亮度变化满足一定假设的情况下具有潜在优势；半直接方法则尝试在计算效率与跟踪精度之间取得折中。对于工业巡检而言，视觉SLAM的主要价值在于能够以相对较低硬件成本获得连续位姿和图像地图，但其适用性高度依赖环境纹理、光照稳定性和动态程度。

激光SLAM的发展则主要围绕高精度几何约束和实时计算展开。LOAM通过区分激光里程计和地图构建过程，在保证实时性的同时获得了较好的定位精度
\cite{zhang2014loam}。随后，激光--惯性融合方法进一步利用IMU提供高频运动先验，使系统能够在快速运动和稀疏点云环境中提高状态估计稳定性。FAST-LIO2直接利用原始点云建立残差，并采用高效的数据结构维护地图，避免了复杂人工特征提取过程
\cite{xu2022fast}；Faster-LIO则进一步围绕轻量化和并行计算降低实时处理开销
\cite{bai2022fasterlio}。对于工业无人机，这类方法尤其适合用于结构规则、空间几何信息丰富而纹理较弱的厂房、管廊和大型设备周围环境，但其计算负担、传感器成本和点云数据存储需求仍需考虑。

视觉--惯性融合具有另一类优势。相机可以提供丰富的环境纹理和语义信息，IMU能够在短时间内提供高频率的运动约束，因此二者结合可以缓解单纯视觉系统在快速运动或短时遮挡下的跟踪不稳定问题。VINS-Mono等系统通过联合优化视觉和惯性测量构建状态估计框架
\cite{qin2018vins}。在微型无人机中，单目相机和IMU组成的轻量感知配置已经能够支持较激进的飞行动作
\cite{loianno2017estimation}，并在自主飞行与实时三维建图中得到验证
\cite{faessler2016autonomous}。然而，对于工业厂区而言，单目视觉的尺度问题、光照变化和重复纹理仍可能导致定位退化，因此工程系统通常需要结合GNSS、激光或其他外部信息补充全局约束。

从融合结构看，紧耦合方法近年来逐渐受到重视。传统松耦合系统通常先由不同传感器独立产生位姿或里程计结果，再进行后端融合，工程实现相对简单，但中间处理可能损失原始测量之间的互补信息；紧耦合方法则在同一个状态估计框架中直接利用原始或低层测量约束，从而能够更充分地利用跨模态信息。对于无人机自主巡检，紧耦合并不意味着一定优于松耦合，而是需要在实时性、计算复杂度、故障隔离能力和工程维护成本之间进行权衡。当系统采用多种传感器并且环境存在频繁退化时，能够独立判断单个传感器是否可信，以及在某一传感器退出后保持系统可运行，往往比单一工况下的极限精度更加重要。

多模态感知还包括跨模态几何与语义融合。激光点云可以提供准确的空间结构，但缺乏丰富的颜色和语义描述；视觉图像具有丰富纹理，却难以直接提供绝对尺度和稳定的三维几何信息。因此，将视觉图像中的设备、部件和缺陷类别映射到点云或三维地图中，可以建立“几何位置+语义类别+异常状态”的统一场景表示。近年来语义SLAM、学习型特征以及主动SLAM逐渐把关注点从单纯定位拓展到“理解环境并决定下一步观测什么”
\cite{huangzx2023survey,placed2023active}。这种变化与巡检任务天然契合：传统SLAM只回答“当前在哪里”，而主动感知开始考虑“去哪里观察可以获得更多有价值的信息”。

针对工业设备检测，视觉和红外的融合尤其具有实际意义。可见光能够识别表面纹理和几何缺陷，红外可以发现温度异常，二者同时存在时能够降低单一模态的盲区。例如，某些电气设备在外观上没有明显变化，但局部温升已经出现；而某些结构裂纹则几乎不存在明显温度差异，仅依赖热像无法可靠识别。多模态系统可以通过图像级、特征级或决策级融合，将结构缺陷与状态异常纳入统一认知框架
\cite{jadin2012recent,usamentiaga2014infrared}. Wang等提出的M3DM进一步在RGB与三维点云之间进行混合融合，并通过多记忆库保留不同模态信息，说明工业异常检测已经从单一图像扩展到RGB-3D联合建模的方向\cite{M3DM2023}。该类思路对于无人机巡检具有启发意义，因为很多工业缺陷不仅表现为二维纹理变化，还可能对应几何变形、凹陷或装配关系变化。

从实际工程案例看，Jeon等提出的面向架空输电设施的自主飞行策略，将光学相机与三维激光信息结合，用于输电塔接近和线路跟踪，并在虚拟与真实环境中进行验证。这类研究将目标检测、三维感知、接近策略和飞行控制组合起来，体现出“多模态感知服务于自主任务”的系统思路，而不再将感知与巡检行为割裂开来。对于复杂工业场景而言，这一方向比单独提升某一感知算法精度更具有工程意义，因为最终任务的评价指标往往是是否安全抵达目标、是否获得有效观测以及是否能够稳定完成巡检，而非某个局部模块在标准数据集上的单一指标。

面向真实巡检任务的研究已经开始尝试把多模态感知直接嵌入自主接近和目标跟踪过程。Jeon等针对架空输电设施提出了结合光学相机与三维激光信息的自主飞行策略，通过多模态信息辅助输电塔接近、线路跟踪以及目标检测，并在虚拟和真实环境中验证了方法的有效性ootnote{该研究的基本信息可由相关文献核验。}\cite{jeon2024multimodalinspection}。该类工作的重要启示在于，多模态融合的最终目标并不是得到一张“更丰富”的感知数据，而是为具体的巡检行为提供可靠约束。例如，视觉可以负责目标识别和纹理提取，激光可以用于估计目标空间位置和相对距离，二者共同决定无人机是否可以安全接近以及相机是否处于合适的观测位姿。由此，感知融合开始从“传感器层融合”进一步向“任务层融合”发展。

此外，多模态融合还应关注模态之间的信息冗余和冲突。若多个传感器提供高度相关的数据，简单增加模态并不一定能够提升系统性能，反而会增加同步、标定和计算开销；若不同模态在某一时刻产生相互矛盾的结果，则系统需要判断冲突来源并选择更可信的证据。因此，工业无人机感知系统未来需要研究基于任务和环境状态的动态模态选择，而非固定配置所有传感器。

复杂工业感知的评价也因此应该从传统的平均定位误差、重投影误差等指标向任务级指标扩展。对于自主巡检系统，需要关注定位是否连续、地图是否可重复、异常观测是否可关联、传感器失效后能否恢复以及在有限算力下能否保持实时性。已有研究在快速自主飞行、GPS拒止和地下受限环境等方面已经验证了多模态感知的有效性
\cite{mohta2018fast,petrlik2020robust}，但工业现场仍需要更多长期运行、跨工况和多传感器退化条件下的实测数据。尤其是对于烟雾、蒸汽、强反光、粉尘与电磁干扰同时存在的环境，单次实验成功并不足以证明系统具有工程鲁棒性。

因此，从研究脉络上看，无人机感知技术已经由“单传感器定位”发展到“多模态紧耦合状态估计”，并进一步向“语义化、主动化和任务驱动感知”演进。未来的关键问题不再只是选择视觉还是激光，而是如何在不同工业环境下动态配置传感器、估计各模态信息的可信度、识别退化状态并使感知结果直接服务于后续异常认知和任务规划。只有建立这种跨模块的信息闭环，感知系统才能真正从自主飞行的附属模块转变为自主巡检的核心基础。



从工程系统设计角度，多模态融合还需要解决传感器空间位置和时间基准不一致的问题。无人机在飞行过程中姿态变化较快，即使相机和激光雷达安装位置固定，不同传感器实际观测到的场景也存在视场差异。如果外参误差、时间同步误差较大，就可能造成图像和点云之间的空间错位，进一步影响三维重建和目标定位。因此，实际工业巡检系统需要把传感器标定和同步作为自主感知的一部分，而不是仅在系统部署前一次完成。长期运行过程中，机体振动、温度变化和载荷更换都可能引起安装关系变化，因此在线标定和自校验也值得研究。

在退化环境下，多模态融合的难点还在于“不同模态可能同时不可靠”。例如强烟雾会影响可见光和部分激光观测，强热源可能造成红外饱和，GNSS遮挡则会使绝对定位约束消失。传统融合方法通常假设至少有一个可靠传感器持续提供有效信息，但工业现场的极端工况可能不满足这一假设。因此，需要研究融合系统的退化模式识别：首先判断当前环境属于何种退化状态，然后选择合适的传感器组合和估计模型。这种由“先估计再融合”向“先判断感知可信度再融合”的思路，更符合复杂工业场景的实际需求。

同时，感知结果还应具有一定的任务可解释性。对于自主巡检而言，地图中某个区域存在高反射点云、视觉特征不足或者定位协方差突然增大，都可能意味着无人机正在进入风险区域。若系统能够把这些信息传递给任务规划器，就可以主动降低飞行速度、增加观测次数或切换传感器，而不是在定位失效之后才进行被动处理。因此，未来多模态感知研究应逐步建立“观测质量--状态估计--任务风险”之间的映射关系。

\section{工业设备缺陷检测与异常认知}

\subsection{基于视觉与多模态信息的工业缺陷检测}

工业设备异常具有明显的层次性和异构性。按照异常形成原因，可以将其分为结构表观异常、部件状态异常、热学异常和环境异常。结构表观异常包括裂纹、锈蚀、剥落、变形和断股等，通常通过可见光图像进行识别；部件状态异常包括部件缺失、松动、破损、异物悬挂等，需要结合设备结构知识判断；热学异常包括接头过热、光伏热斑和绝缘劣化发热等，通常更适合通过红外测温发现；环境异常则可能表现为施工机械入侵、烟火、通道障碍和人员进入禁区等
\cite{taox2021survey,nguyen2018automatic,miao2020uav}。不同类型异常的成像规律和诊断依据并不相同，这意味着工业巡检不能简单地使用一个视觉分类器解决全部问题，而需要依据任务和异常机理选择合适的观测模态与认知方法。

从检测粒度看，工业异常认知至少可以分为图像级分类、目标级检测和像素级分割三个层次。图像级分类主要判断当前观测是否存在异常，适合用于快速筛查；目标检测进一步给出异常对象的位置，适合用于设备部件级定位和巡检记录生成；像素级分割则可以精细描述裂纹、腐蚀和剥落等缺陷的形态，为后续面积、长度或严重程度计算提供依据
\cite{taox2021survey}. 在无人机巡检中，三种任务往往是串联的：无人机首先利用轻量模型进行快速筛选，当发现疑似异常后进一步靠近目标，再采用高分辨率模型或分割模型进行精细分析。这种“粗筛+精检”的分层认知模式比单一模型在所有阶段都采用高复杂度推理更加符合机载资源受限的实际情况。

传统计算机视觉通常依靠人工设计的边缘、纹理、颜色和形状特征进行缺陷识别，但工业现场复杂背景和缺陷尺度变化很快限制了其泛化能力。深度卷积网络的出现使模型能够从数据中自动学习层次化特征，显著推动了工业缺陷检测的发展。AlexNet开启了大规模深度视觉学习在图像识别领域的应用，VGG、ResNet等网络进一步提升了特征表达能力
\cite{krizhevsky2012imagenet,simonyan2015very,he2016deep}。随着网络结构不断成熟，工业检测逐渐从分类任务扩展到目标检测和像素级分割。Faster R-CNN等两阶段方法通过区域候选与精细回归获得较高检测精度，而YOLO系列、SSD和RetinaNet等单阶段方法更加关注实时性
\cite{ren2015faster,redmon2016you,liu2016ssd,lin2017focal}. 对于无人机在线巡检而言，两阶段方法更适合在地面服务器或高算力平台上进行精检，而轻量单阶段模型更适合机载实时任务。

近年来，Transformer结构逐步进入工业视觉检测领域。DETR将目标检测转化为端到端集合预测问题，Deformable DETR进一步通过稀疏注意力机制提高了对多尺度目标的处理能力
\cite{carion2020end,zhu2021deformable}。视觉Transformer和Swin Transformer利用自注意力建模全局关系，对复杂背景下的目标表征提供了新的解决思路
\cite{dosovitskiy2021image,liu2021swin,vaswani2017attention}。但从工业巡检角度看，模型性能提升不能只用标准数据集上的平均精度衡量。无人机采集的图像具有明显的距离、视角和姿态变化，缺陷常常只占据很少像素，因此小目标召回率、近距离推理延迟以及不同高度下的稳定性更加关键。

针对工业巡检中的小目标问题，研究通常从多尺度特征、特征融合、注意力机制和高分辨率输入等方向改进检测网络。输电线路中的销钉、螺栓、绝缘子部件以及电线缺陷往往具有较小的图像占比，稍微的运动模糊和压缩失真就可能导致漏检
\cite{nguyen2018automatic,tao2020detection,zhouwq2024review}。桥梁裂缝等细长缺陷也具有类似问题，其宽度通常远小于背景结构尺度，因此单纯依靠常规目标检测框进行定位会造成边界信息损失，需要进一步利用分割或像素级表征
\cite{yang2018pixel,yeum2015vision,cha2017deep}. 在无人机实际系统中，还应根据飞行高度和目标尺寸动态决定图像分辨率、相机焦段和推理策略，使“飞行控制”和“检测性能”形成联动。

深度迁移学习在一定程度上缓解了工业缺陷数据不足的问题。对于新设备或新厂区，可以利用在大规模数据集上预训练的视觉特征作为基础，再通过少量目标领域数据进行微调
\cite{gao2018deep}. 这种策略能够降低从零开始训练模型的成本，但其本质仍然依赖目标场景中存在一定数量的代表性样本。当设备型号发生变化、拍摄角度发生变化或正常外观本身具有较大波动时，迁移后的模型可能出现性能下降。因此，工业巡检中的泛化问题与通用视觉分类存在明显差别：模型不仅要适应新类别，还要适应新设备、新工况、新环境和新的成像条件。

公开数据集的建立推动了工业缺陷检测研究的标准化。MVTec AD等数据集将工业异常检测从单一项目验证扩展到统一基准评测，并推动了无监督异常检测的发展
\cite{bergmann2019mvtec}。InsPLAD等无人机电力设施数据集进一步将真实航拍视角、部件类别和缺陷标注纳入公开基准，为无人机巡检算法评测提供了更加接近实际的数据来源
\cite{vieira2023insplad}. 但是，公开数据集与真实工业现场之间依然存在域差异。实验室数据往往具有稳定的光照、固定的相机位置和较为完整的目标，而无人机现场数据存在运动模糊、曝光变化、距离变化、压缩损失和环境遮挡。因此，在评价无人机巡检算法时，应同时关注标准数据集性能和真实飞行条件下的表现。

红外热像的引入使工业异常认知从“外观检测”扩展到“运行状态检测”。电气设备的局部过热、轴承温升和光伏组件热斑等异常往往在可见光图像上不明显，但通过温度场能够直接获得异常线索
\cite{jadin2012recent,usamentiaga2014infrared}。然而，红外数据本身也受环境温度、反射、发射率以及距离影响，因此简单地采用固定温度阈值容易造成误报或漏报。未来更合理的做法是将红外与可见光、设备运行参数以及历史状态联合分析，从“温度异常”进一步推断“设备异常”。这意味着工业巡检的多模态融合不应止步于图像拼接或特征连接，而应进一步利用不同模态的物理意义和时空对应关系。

在多模态异常检测方面，RGB-3D联合建模提供了另一个值得关注的方向。M3DM通过RGB特征和点云特征的混合融合，利用跨模态对齐和多记忆库保留不同模态的信息，对MVTec-3D AD等任务展示了多模态融合的有效性。该思路对于无人机巡检的意义在于，某些设备异常同时表现为表面纹理变化和几何变化，例如凹陷、凸起、结构偏移等，仅利用二维图像难以充分刻画。将图像、点云和历史三维模型关联后，可以进一步开展三维缺陷定位和变化检测，使异常结果从“某张图片里有问题”变为“设备三维模型中的某一部位发生了变化”。

总体来看，监督式工业缺陷检测已经能够较好地解决“已知设备+已知异常类型”的问题，且通过轻量化和多尺度建模已经具备一定的工程部署基础。但复杂工业无人机巡检还面临三个更加突出的限制：一是缺陷样本数量少且极不均衡；二是飞行视角、成像环境和设备型号变化造成明显的域偏移；三是实际巡检并不只需要一次性检测，还需要给出空间位置、异常等级和历史变化。由此，研究重点正在从“如何提高已知缺陷的检测精度”逐步转向“如何在缺陷样本不足和异常未知的情况下可靠判断设备状态”。



对于无人机巡检，检测模型的输入数据具有明显的时空连续性。传统工业视觉检测大多处理固定工位摄像机采集的图像，同一设备在相邻帧中的位置变化较小；无人机则会不断改变视角、距离和姿态，因此同一缺陷可能在不同帧中呈现完全不同的尺度和背景。单帧检测如果独立进行，容易出现同一异常多次报警、短时漏检或置信度剧烈波动。因而，在实际巡检中需要进一步利用视频时序信息和目标跟踪，将多个观测帧融合为一个稳定的异常事件。这样既可以减少重复报警，也能够通过多视角证据提高异常判断置信度。

另一个问题是缺陷检测与无人机飞行距离之间存在直接耦合关系。当无人机距离设备较远时，只能得到粗粒度目标信息，此时适合进行快速筛查；当检测到疑似异常后，系统应主动缩短与目标的距离，以获得更高空间分辨率。这意味着目标检测模型实际上可以成为路径规划的触发器：检测到高风险目标时输出其空间位置和置信度，由规划模块计算最佳复检视点；当多次观测仍无法确认时，再根据异常类型选择红外、高清变焦或三维扫描等不同传感器。因此，工业无人机缺陷检测并不是独立的软件模块，而是主动巡检闭环中的一环。

2025年的相关综述进一步表明，工业视觉异常检测已经形成监督、半监督、弱监督、自监督和无监督等多种学习范式，研究重点开始从单纯追求检测精度扩展到数据获取、模型泛化和实际部署等问题\cite{li2025iad}。与此同时，针对RGB、3D及多模态无监督工业异常检测的研究表明，RGB数据具有丰富的颜色和纹理信息，3D数据能够补充空间几何信息，而多模态方法虽然能够降低单一模态的信息缺失，却同时带来了跨模态配准、特征干扰和计算成本增加等新问题\cite{lin2025rgb3dmm}。这些结论与无人机巡检场景具有较强的一致性：无人机采集的数据天然具有多视角、三维和时序特征，因此异常检测方法需要在“信息更丰富”和“系统更复杂”之间进行权衡，而不能简单地以增加传感器和模型规模作为性能提升的唯一途径。

从评价指标看，工业巡检也不宜只使用mAP或IoU等通用指标。对于现场应用，漏检一个关键设备缺陷的代价可能远高于多报几个普通异常，因此召回率、虚警率、检测时延以及异常发现后的复检成功率更加重要。对于细小裂纹、腐蚀点和螺栓缺失等任务，还需要关注目标尺寸与成像距离之间的关系，并通过实际飞行实验建立“距离--分辨率--检测性能”的映射。只有把模型性能与飞行任务联系起来，才能判断一个检测方法是否真正适合无人机巡检。

\subsection{小样本、无监督异常检测与基础模型认知}

工业异常检测具有典型的小样本和长尾特征。正常设备在绝大多数时间都处于正常状态，而严重缺陷往往低频出现，不可能为了训练模型而人为大量制造真实故障。与此同时，不同设备的正常外观和运行特征也存在差异，使得一个统一的监督式分类模型难以覆盖全部场景。因此，在工业领域，“正常样本比异常样本容易获得”成为推动无监督和自监督异常检测发展的重要现实条件
\cite{taox2021survey,pang2021deep}. 这与传统目标检测的假设明显不同：传统检测希望学习“目标类别是什么”，而工业异常检测更希望学习“什么是正常”，再把偏离正常状态的样本识别为异常。

早期无监督异常检测主要采用重建和生成式建模思路。VAE通过概率生成模型学习正常数据分布，GAN则通过生成器和判别器学习正常图像结构
\cite{kingma2014auto,goodfellow2014generative}。f-AnoGAN进一步利用GAN隐空间实现异常定位，而Deep SVDD通过学习正常样本的紧致特征边界来进行单类判别
\cite{fangan2019,ruff2018deep}。这些方法的重要意义并不只是提出了新的网络，而是改变了工业异常检测的数据逻辑：模型训练阶段不必等待大量真实缺陷出现，而可以直接利用数量庞大的正常样本建立正常状态表征。

随后，基于预训练特征和记忆库的方法成为工业异常检测的重要路线。SPADE通过多尺度特征对应进行异常判断，PaDiM通过对局部特征分布进行概率建模，PatchCore则建立正常特征记忆库并通过核心集采样降低计算开销
\cite{roth2022towards}. 此类方法的一个突出特点是，不需要针对每一种新缺陷重新设计分类器，而是通过正常样本特征空间中的偏离程度进行异常定位。对于无人机巡检而言，这种思路很有吸引力，因为现场通常容易采集大量正常设备图像，而难以获得完整的缺陷类型集合。然而，正常样本本身并非固定不变。不同光照、不同角度、不同设备状态都可能被算法误认为异常，因此工业无人机应用必须把“环境变化”与“设备异常”区分开来。

自监督异常检测进一步降低了对人工异常样本的依赖。CutPaste通过将正常图像局部剪切并重新粘贴制造伪异常，从而训练模型学习局部变化；DRAEM通过重建表征与判别分割的组合提升异常定位能力
\cite{zavrtanik2021draem}. 这类方法的核心思想是利用人工构造的监督信号，使模型在只有正常数据的条件下获得一定的异常敏感性。但人工合成的异常与真实工业缺陷之间存在差异，如果模型过度依赖合成伪影，可能在真实异常出现时出现泛化下降。因此，异常合成的真实性、可控性以及与真实缺陷机理的一致性，是自监督工业异常检测值得长期研究的问题。

EfficientAD等研究进一步强调了工业在线检测对实时性的要求。高精度模型如果只能在离线服务器上运行，就难以直接支撑无人机机载实时筛查。对于无人机而言，算法需要在有限计算功耗和飞行续航条件下完成持续推理，因此“精度+速度+能耗”的综合指标比单一精度指标更加重要。另一方面，机载系统通常更适合完成疑似异常的快速发现，而地面边缘服务器或云平台则可以对高分辨率数据实施复杂的二次分析，因此未来工业巡检很可能采用分层计算架构，将快速筛查和精细分析分别部署在不同计算层。

近年来，视觉语言模型和视觉基础模型逐渐进入工业异常检测领域，为小样本甚至零样本识别提供了新的路径。CLIP通过大规模图像--文本对比学习获得了较强的跨类别表征能力
\cite{radford2021learning}。WinCLIP将CLIP进一步用于工业异常分类与分割，通过窗口特征与文本提示建立正常/异常语义关系，在MVTec AD和VisA等基准上展示了零样本和少样本异常检测的潜力。与传统方法相比，其最大价值在于减少对每一个工业类别单独训练模型的依赖，使“一个通用视觉模型适应多个巡检对象”成为可能。

基础模型的另一项重要价值是开放词汇目标理解。SAM提供了通用提示分割能力，可以根据点、框等提示对复杂目标进行分割
\cite{kirillov2023segment}；Grounding DINO通过文本与图像之间的开放词汇关联支持目标检测
\cite{liu2024grounding}；DINOv2则利用大规模自监督学习获得具有迁移能力的视觉特征
\cite{oquab2024dinov2}。这些技术为工业巡检提供了“目标先验+异常检测”的新组合形式。例如，对于一个新建工业装置，即使设备缺陷没有被完全标注，系统仍可以利用文本或结构知识定位“阀门”“法兰”“绝缘子”等巡检对象，再在对象内部进行异常检测。这与传统“先建立固定类别数据集，再训练专用检测器”的方式存在明显差别。

从基础模型在异常检测中的作用来看，近期研究开始将其划分为特征编码器、异常检测器和结果解释器等不同角色。相关综述指出，基础模型的主要优势在于能够利用大规模预训练获得较强的通用表征，并在少样本、零样本以及跨模态任务中展现潜力；但同时存在模型计算量大、提示敏感、领域偏差和解释可靠性不足等问题\cite{ren2025foundationAD}。另一项面向工业缺陷检测的综述也指出，基础模型尤其适合标注数据稀缺的场景，但模型规模带来的推理速度下降使其在实时工业应用中面临部署挑战\cite{yang2025foundationIAD}。因此，基础模型与轻量化检测器、规则知识和多传感器证据相结合，更可能成为无人机工业巡检中的实际技术路线。

AnomalyGPT等研究进一步尝试利用大视觉语言模型直接进行工业异常对话和解释，使系统能够不仅输出异常分数，还能够对异常位置进行描述并支持多轮交互\cite{gu2024anomalyGPT}。对于工业巡检而言，这种能力的意义在于，最终使用者通常并不只关心“异常概率是多少”，而更关心“哪个设备出了什么问题、严重程度如何、是否需要复检”。然而，大模型的开放能力也带来了新的问题：模型可能产生与视觉证据不完全一致的解释，异常判断的阈值和不确定性量化也更加复杂。因此，在安全关键的工业场景中，基础模型更适合作为认知增强模块，而不能直接替代基于物理和视觉证据的确定性诊断。

PromptAD进一步探索了只利用正常样本学习异常提示的方法，说明提示学习已经开始从手工设计逐步走向数据驱动\cite{promptad2024}。RealNet则将扩散模型用于构造更加逼真的合成异常，并通过特征选择降低冗余\cite{realnet2024}。SimpleNet通过预训练特征适配、特征空间异常生成和轻量判别器构建了一套较为简洁的工业异常检测体系
\cite{promptad2024,realnet2024,simpleNet2023}. 这些方法从不同角度说明，工业异常检测正在形成“预训练特征+正常样本+少量先验”的新范式\cite{simpleNet2023,realnet2024,promptad2024}。对于无人机巡检，这一范式比大规模标注监督学习更接近真实工程条件，因为现场新增设备或新增缺陷通常缺乏充分的异常样本。

值得关注的是，多模态异常检测正在由二维RGB向RGB、3D和热成像等多源信息扩展。RGB图像更擅长表现表面纹理，点云更擅长表现几何变化，红外则能够反映设备热状态。当多个模态能够在统一空间和时间基准下对应到同一设备时，就有可能通过模态间的一致性和互补性提高异常判断的可靠性。例如，如果图像显示结构外观发生变化，同时三维点云也显示局部几何偏移，那么异常的可信度往往高于仅依据单张图像得到的结果；相反，如果某个异常只出现在一幅图像中而在历史三维模型和其他模态中均不存在，系统可以将其判定为低置信度结果并触发复核。

从“异常检测”进一步走向“智能预警”，还需要时间维度。单次巡检只能告诉系统当前是否存在异常，而无法回答异常是否在持续发展。真正的工业状态评估需要把多次飞行获得的图像、点云和温度信息进行时空对齐，建立设备异常的历史轨迹。对于长期运行的设备，可以比较同一部件在不同日期、不同工况下的视觉特征、几何形态和热状态变化，从而形成趋势性指标。这一过程又与数字孪生和工业数据平台形成自然联系：无人机负责获取现场变化，数字孪生负责组织设备对象和历史状态，异常认知模块负责将原始数据转化为可解释的设备健康信息
\cite{tao2019digital}。

近年来的工业异常检测综述也强调了多种监督范式、数据模态和真实部署条件之间的差异。已有工作指出，工业异常检测不仅需要比较算法在标准数据集上的AUROC等指标，还需要考虑数据采集、数据噪声、标签质量、实时推理以及不同工业场景中的域偏移。2025年的相关综述进一步系统梳理了监督、半监督、自监督、弱监督和无监督路线，显示工业异常检测已经从单一算法问题逐渐发展为涵盖数据、模型和部署全过程的系统问题。基于基础模型的综述则进一步指出，基础模型虽然特别适合少样本和零样本场景，但模型规模、推理延迟和工业领域适应性也是实际部署中的重要矛盾。

因此，从研究发展脉络看，工业异常认知正在经历“监督检测--小样本学习--无监督异常检测--开放世界认知”的演进。监督学习在已知缺陷识别方面仍具有较高的精度优势，无监督和自监督方法降低了缺陷标注成本，视觉语言模型和基础模型则使模型开始具备跨类别、开放词汇和少样本能力。但对于复杂工业无人机巡检，仍不能简单地把公开基准上的高性能视为工程问题已经解决。一方面，真实现场存在明显的环境域差异；另一方面，未知异常可能与正常工况变化具有相似外观。未来更具价值的方向，是把视觉和热成像、三维结构、历史记录以及设备知识联合起来，在保持检测性能的同时提供可解释的不确定性，并把认知结果进一步反馈给巡检规划。



小样本异常检测的核心难点还在于正常样本本身具有多样性。在一个真实工业设备中，温度、季节、运行负载、表面污染和光照变化都可能使正常图像产生较大差异。如果模型把这些正常变化压缩到过于狭窄的“正常分布”中，就会提高误报率。因此，未来工业异常检测需要从“学习正常外观”进一步转向“学习正常工况范围”。对于无人机而言，可以利用飞行时间、设备运行状态、环境温度以及相机姿态等上下文变量对视觉特征进行条件化建模，使模型判断“在当前工况下是否异常”，而不是只判断“与训练图片是否相似”。

开放世界异常认知还需要解决异常定义的问题。工业现场的异常并不总是一个明确的类别，有些异常属于新型故障，有些只是设备工艺状态发生变化，还有一些可能是传感器自身造成的伪异常。因此，未来基础模型需要具备区分“设备异常、环境变化和传感器异常”的能力。视觉语言模型能够通过文本提示表达较高层次语义，但其输出仍需要由图像、红外、三维数据和设备知识共同验证。比较合理的系统结构可能是以基础模型承担开放语义理解，以专用检测器承担细粒度定位，再通过规则或概率模型完成最终风险判定。

异常预警进一步涉及风险量化。如果系统只输出正常/异常二值结果，就无法支撑巡检任务排序。实际应用更需要连续性的风险指标，例如异常置信度、严重程度、发展速度以及潜在后果。可将单次检测结果与历史巡检数据融合，形成设备级健康指数，再根据阈值触发复检、人工确认或维护工单。这样，异常检测才真正从视觉任务转变为工业状态认知任务。

\section{无人机自主巡检规划与多无人机协同}

\subsection{自主巡检任务规划与动态路径重规划}

无人机自主巡检任务规划与普通无人机导航存在本质区别。普通导航的主要目标是从起点安全到达终点，而巡检任务需要在安全飞行的前提下对设备表面或空间区域进行充分观测，因此不仅关注路径长度，还关心目标覆盖率、观测角度、距离约束和图像质量。对于具有明确结构的工业设备，还需要针对设备表面制定一组合理的观察位姿，使相机视场能够覆盖关键区域，同时避免无人机过度接近障碍物。由此，工业巡检规划本质上是一个“运动可行性+信息获取价值+任务优先级”的联合优化问题。

经典路径规划方法可以从图搜索、采样搜索和局部优化等角度理解。Dijkstra和A*利用离散图搜索获得从起点到终点的可行路径，结构清晰、易于解释，但依赖于环境地图并容易受到栅格分辨率影响
\cite{dijkstra1959note,hart1968formal}。PRM与RRT适合高维连续空间，能够在复杂构型空间中搜索可行路径；RRT*进一步通过重布线提高路径的渐近最优性
\cite{kavraki1996probabilistic,lavalle2001randomized,karaman2011sampling}. 这些方法为无人机全局路径规划建立了成熟基础，但原始形式通常没有针对巡检任务设计，也没有充分考虑四旋翼的动力学约束和感知价值。

在动态环境中，环境信息会随着时间变化，传统一次性规划得到的路径可能很快失效。D*和D* Lite通过增量式搜索机制，在地图局部变化后只更新受影响区域，因此比完全重新规划更适合在线调整
\cite{stentz1994optimal,koenig2002d}. 局部避障方法则更加关注短时间范围内的安全轨迹，如人工势场、动态窗口和互惠速度障碍等方法能够根据局部障碍物位置快速产生避障动作
\cite{khatib1986real,fox1997dynamic,van2011reciprocal}. 但对于工业无人机而言，仅根据几何障碍物进行局部避障还不够，因为设备本身既是障碍物，又是需要被观测的目标。因此系统必须判断“应该绕开设备还是接近设备”，并根据巡检阶段动态改变行为模式。

覆盖路径规划（Coverage Path Planning, CPP）为大范围设施巡检提供了更直接的理论基础。Choset和Galceran等系统总结了覆盖规划的基本概念和环境分解方法
\cite{choset2001coverage,galceran2013survey}. 无人机覆盖规划通常需要将作业区域分解为多个可覆盖区域，再根据传感器视场、航迹宽度和转弯代价设计路径。对于厂区巡检，除了二维平面覆盖之外，还需要处理立体设备和垂直表面。例如储罐外壁、炉体、塔架和管廊不能简单地投影为地面区域，而需要在三维空间中定义观察位姿。由此，三维覆盖规划与视点规划逐渐成为设施巡检中的重要问题。

下一最佳视点（Next-Best-View, NBV）方法强调利用当前地图和不确定性主动选择下一观察位置。Bircher等的工作表明，若无人机能够根据尚未观测区域和信息增益选择下一视点，就可以比盲目执行固定航线获得更高的信息效率
\cite{bircher2016receding}. 对于工业巡检而言，这一思想尤其重要，因为巡检并不一定需要均匀覆盖所有空间。当系统已经判断某设备处于正常状态时，可以降低进一步观测频率；当发现疑似缺陷时，则可以主动规划更近距离、多角度甚至重复观测，从而提高诊断置信度。此时，规划目标由传统的“最短距离”转变为“单位飞行资源获得最大有效信息”。

随着轨迹优化的发展，规划结果越来越重视无人机动力学可行性。Mellinger和Kumar提出的最小snap方法将位置、速度、加速度和高阶导数纳入轨迹生成，使四旋翼能够在动态约束下执行高速飞行
\cite{mellinger2011minimum}。Usenko等利用均匀B样条和局部地图结构实现了实时轨迹重规划，EGO-Planner进一步降低了复杂ESDF地图依赖并提高了快速自主飞行中的局部轨迹生成效率
\cite{usenko2017real,zhou2022ego}. 这些方法对于工业巡检的重要启示是，路径规划不能脱离飞行控制独立存在。狭窄管廊或复杂钢结构环境中的可行航迹不仅需要避障，还必须满足速度、加速度、角速度和安全距离等约束，否则即使几何路径可行，也可能由于控制执行困难导致实际任务失败。

近年来，面向具体基础设施的研究开始把任务分配、轨迹生成和任务时序放在同一个规划框架中。以风电机组巡检为例，Silano等采用信号时序逻辑描述任务执行关系，并结合混合整数规划与轨迹约束实现多飞行器协同任务与轨迹规划，研究不仅关注是否覆盖巡检区域，还关注任务发生的时序和无人机之间的安全关系，并通过仿真与真实试验进行验证	extsuperscript{}。这说明巡检规划正逐渐从传统的“几何路径优化”向“任务逻辑+轨迹安全+多机协同”的联合规划转变。对于大型工业厂区而言，同类思想可以进一步扩展到常规巡检、临时巡检和异常复检之间的任务切换。

目前工业巡检系统中比较常见的方式仍然是预先建立航线或者巡检点序列，在飞行过程中只进行局部微调。这种方法工程实现简单，特别适合设备结构稳定、巡检任务重复的场景。但其局限性也很明显：它假设环境和任务需求基本不变，而真正的工业现场往往会不断产生新信息。例如一次巡检发现设备异常后，原先设定的路径可能需要立即增加复检点；某个作业区域临时出现人员或车辆时，需要重新规划；无人机剩余电量下降时，还需要改变目标顺序。由此，未来自主巡检规划需要由“航线执行”逐渐转向“状态驱动的任务规划”。

状态驱动规划意味着感知和规划之间必须存在双向信息流。传统链路是“规划产生航线，感知执行航线”，而闭环巡检应该是“规划产生初始任务--感知获得设备状态--异常认知改变任务优先级--规划重新决定观察位姿--再次感知确认状态”。这一闭环与主动感知、下一最佳视点以及信息增益规划天然相连。其难点在于异常认知输出往往具有不确定性，而规划又需要在实时性与信息价值之间作出权衡。例如当模型给出一个中等置信度的腐蚀异常时，系统究竟应该立即靠近复查，还是先完成附近其他高优先级任务，需要综合考虑风险等级、复查成本和剩余能量。

因此，工业巡检规划的研究正在从“几何可行路径”逐渐转向“任务语义约束下的自主行为”。未来规划模型需要同时考虑设备重要程度、异常严重程度、图像质量需求、三维可见性、剩余电量和通信状态，并能够处理连续任务和临时任务之间的动态切换。与单纯追求最短路径相比，这类规划更接近实际工业作业过程，也更能体现“自主巡检”而不是“自动飞行”的本质区别。



对于工业巡检，任务规划还可以按照任务粒度分为区域级、设备级和部件级三个层次。区域级规划解决无人机如何覆盖厂区或装置区的问题，设备级规划解决如何围绕储罐、塔架、风机和输电塔等目标建立观测序列，部件级规划则解决如何选择视角和距离以获取足够分辨率的局部数据。三者之间存在明显的层级关系：区域级规划决定访问顺序，设备级规划决定观测路径，部件级规划则根据异常情况执行局部复检。当前许多路径规划研究主要针对其中某一层，而实际自主巡检需要将三个层次结合起来。

巡检任务还存在典型的“先粗后细”特点。第一次飞行可以采用较高速度完成设备普查，目标只是获得每个设备的基本状态；当识别出疑似异常后，再自动生成局部精检任务。与一次性将全部设备按照最高精度巡检相比，这种自适应策略可以将有限能源集中到高价值目标上。在信息论意义上，其目标不是最大化覆盖面积，而是最大化单位能量和单位时间内的有效状态信息获取量。因此，未来路径规划可以将异常概率、信息增益和复检收益直接引入代价函数。

动态重规划还需要处理安全约束的优先级。当系统同时面临“高风险设备需要复检”和“附近出现动态人员”两个事件时，规划器不能单纯追求任务收益，而必须首先满足安全约束。由此，工业自主巡检规划更适合采用分层约束结构：硬约束负责禁飞区、最小安全距离和最低剩余电量等不可违反条件，软约束负责任务收益、观测质量和路径长度等优化目标。将安全约束与任务价值分离，有助于提高规划器在复杂工业环境中的可验证性。

此外，巡检任务不是一次飞行结束后才开始下一轮规划，而是持续运行的过程。无人机在飞行中不断产生新的环境地图和设备状态信息，规划器应持续更新任务优先级。因此，可将自主巡检理解为一个滚动时域决策问题：在当前时刻根据已有信息规划未来一段时间的动作，仅执行其中一部分，然后根据新获取的数据再次更新计划。这种思想与Receding Horizon和主动感知方法一致，也更适合动态工业环境。

\subsection{多无人机任务分配与协同决策}

当厂区范围较大、设备点位较多或需要同时开展常规巡检与临时复检时，单架无人机容易受到续航和作业周期的限制。多无人机协同能够通过空间分工、任务并行和故障备份提高整体巡检效率，因此成为大规模自主巡检的重要技术方向。与一般多机器人任务分配相比，多无人机巡检需要特别关注空间覆盖、任务时序、能源消耗、通信连接和观测资源等问题
\cite{gerkey2004formal,jiagw2021review,duyh2020survey,jiayn2020recent}。

任务分配是多无人机协同的基础。市场机制和拍卖机制通过定义任务收益、飞行代价和资源消耗，让不同无人机根据自身状态竞标任务；CBBA利用一致性协议实现分布式任务分配，在不依赖中央调度器的条件下维持多机对任务归属的基本一致性
\cite{choi2009consensus}. 这种方法具有可解释、易于融合约束等优点，适合任务目标明确且无人机能力差异有限的场景。国内研究还进一步引入粒子群、博弈论和混合优化等方法处理多目标、异构和动态任务分配问题
\cite{zhaiz2023market,zhangrp2022hybrid,jukai2022task,biwh2024review,zhangkw2021survey}。

在静态任务条件下，任务分配可以在起飞前一次性完成；然而真实巡检任务往往存在较强的动态性。某架无人机可能因为低电量提前返航，某个设备可能在巡检过程中被识别为高风险目标，新出现的异常又可能产生临时任务，因此任务分配需要支持在线重分配。在线重分配不是简单地重新运行一次优化算法，因为重新分配本身也会产生通信、计算和飞行成本，并可能造成已有航迹频繁改变。对于高可靠工业巡检，更合理的策略是设置任务优先级和重规划触发条件，在异常等级超过阈值、无人机状态发生明显变化或通信链路发生故障时再执行任务迁移。

多无人机协同的底层理论来自一致性控制和群体行为控制。Vicsek模型和Reynolds模型揭示了局部交互作用如何形成群体有序行为
\cite{vicsek1995novel,reynolds1987flocks}；Olfati-Saber和Ren等则从一致性理论角度研究多智能体系统在动态拓扑中的协同行为
\cite{olfati2004consensus,ren2005consensus,olfati2006flocking}。这些理论为无人机保持队形、同步状态和分布式决策提供了基础。然而，工业巡检中的多机协同并不要求所有无人机一直保持固定队形，更强调“按任务需要形成动态协作关系”。例如，两架无人机可以分区执行常规巡检，当某一区域发现异常时，附近无人机可以临时形成协同小组，一架负责保持安全观察位置，另一架负责近距离复检。

近年来，多智能体强化学习（Multi-Agent Reinforcement Learning, MARL）为复杂协同决策提供了数据驱动的新方法。MADDPG、QMIX和MAPPO等方法从不同角度解决多智能体合作、信用分配和策略学习问题
\cite{lowe2017multi,rashid2018qmix,yu2022surprising,gronauer2022marl}。对于无人机巡检而言，多智能体强化学习的优势在于可以直接学习任务状态、邻机状态和环境信息之间的策略关系，尤其适合任务规则较难显式建模的场景。但其工程应用仍受到训练成本、仿真到现实迁移以及安全性约束影响。工业现场无法像游戏环境那样允许智能体通过大量试错获得策略，因此实际部署通常需要把规则约束、模型预测和学习策略结合起来，而不是完全采用黑盒强化学习。

多无人机协同进一步受到通信条件的限制。工业厂区建筑密集、钢结构众多，通信链路可能随无人机位置不断变化；当无人机飞入地下或设备背面时，甚至可能出现临时失联。传统集中式任务分配假设中央节点能够持续获得所有无人机和任务状态，这在通信稳定时具有较高效率，但通信异常时容易产生单点故障。去中心化方法利用局部感知和邻机信息，使系统即使与中央节点暂时失联也能够继续完成部分任务，因此在复杂工业环境具有潜在优势
\cite{zeng2016wireless,mozaffari2019tutorial}。但完全去中心化又会带来一致性维护和全局最优性下降的问题，如何在集中式与分布式之间取得平衡，是工程系统需要考虑的重要问题。

能源和计算约束同样会显著影响任务分配。无人机不仅要飞到目标区域，还要携带相机、红外或激光设备并持续进行感知计算，因此不同任务的能耗并不只与距离相关。高速飞行、悬停、频繁加减速以及高负载图像处理都会产生额外能源消耗。已有研究从能量高效通信与轨迹优化角度研究无人机资源分配
\cite{zhang2019energy}，也有研究讨论了无人机与移动边缘计算结合以提高计算效率
\cite{zhou2018computation}。将这些研究引入巡检场景后，可以进一步形成“任务收益--飞行能耗--计算代价--通信代价”的综合任务分配模型。

多无人机协同还需要考虑异构平台问题。实际系统中的无人机可能搭载不同焦段相机、热像仪或激光雷达，其飞行速度、续航时间和观测能力各不相同。小型无人机适合进入狭窄区域近距离观察，大型平台适合承担长时间巡检和大范围航线任务。因此，任务分配不能只把无人机视为完全相同的执行主体，而应根据任务需求选择最匹配的平台。对于异常复检任务，还可以按照“发现无人机”和“确认无人机”分工：前者快速巡检并发现疑似异常，后者携带高分辨率或多模态载荷进行精细确认。

近年来真实飞行实验表明，多无人机系统已经逐步从仿真进入实测阶段。相关基础设施巡检研究还将任务分配、轨迹约束和时序逻辑结合到统一框架中，并在风电机组场景进行仿真与真实试验\cite{silano2024stl}。Zhou等展示了微型飞行机器人集群在野外环境中的自主飞行能力，并进一步提出去中心化时空轨迹规划和RACER等协同探索系统，实现依赖局部信息和局部通信的自主协同
\cite{zhou2021decentralized,zhou2022swarm,zhou2023racer}。针对基础设施巡检，Silano等进一步研究了多飞行器任务和轨迹协同问题，将任务分配、轨迹规划以及时序逻辑结合起来，并通过风电机组巡检案例进行验证。这些研究表明，多机系统正在由“多架无人机同时飞”向“多架无人机共同完成一个任务”演进。

但是，从真实工业巡检角度看，目前多无人机研究仍存在明显不足。首先，大量算法验证依赖仿真环境，而真实厂区的通信、动态障碍和安全风险更复杂；其次，多机协同通常优化路径长度、覆盖率或任务完成时间，却较少直接把设备异常等级纳入决策；再次，异构载荷带来的观测能力差异尚未被充分考虑；最后，感知和异常检测结果往往没有实时进入任务分配器。因此，未来多无人机协同应从“空间任务并行”进一步走向“状态驱动协同”，即由设备风险和巡检结果决定谁去哪里、什么时候去以及需要执行何种观测任务。

综合来看，多无人机巡检已经形成从集中式任务分配、分布式拍卖到多智能体学习协同的发展路线，实际飞行实验也在不断增多。但其真正进入工业现场仍需要解决通信受限、异构平台、在线重分配和安全验证等问题。更重要的是，多机协同不应独立于异常认知模块存在，而应与前面的感知和认知形成闭环：一个设备的异常不仅是一个检测结果，同时也是一个高优先级任务事件；一个无人机的能源和通信状态也不仅是资源信息，同时会反过来影响异常复查能否及时完成。由此，自主巡检最终需要从“无人机群规划”走向“面向设备状态的群体智能决策”。



在多无人机巡检中，还应区分“任务分配”和“协同执行”两个层次。任务分配决定哪架无人机负责哪个区域或设备，协同执行则决定多架无人机在执行过程中如何保持必要的信息共享和安全关系。例如，当两架无人机分别负责相邻区域时，通常不需要持续共享全部传感器数据，只需交换位置、速度、任务状态和冲突信息即可；当其中一架发现高风险异常并请求另一架进行复检时，则需要临时共享目标位置、异常类别和最佳观察位姿。因此，通信资源也应作为任务协同的一部分进行动态分配。

工业现场的多无人机协同还具有“任务异步性”。不同无人机的飞行速度、剩余电量、载荷和任务完成时间并不一致，因此难以假设所有智能体在同一时间步同步决策。真实系统更接近事件驱动协同：某无人机完成任务、发现异常、低电量告警或通信状态变化时触发一次局部任务重分配。相比固定周期同步决策，事件驱动方式能够减少不必要的通信和计算开销，但需要解决异步状态下的一致性问题。

多智能体强化学习在此类场景中的潜力主要体现在策略适应性，而不是完全替代传统任务规划。工业巡检可以把规则和优化方法负责“安全和基本可行”，把学习模型用于“复杂情形下的策略选择”。例如，当多个设备同时发生疑似异常时，可以利用学习策略估计不同任务序列的长期收益，再由规则层检查能源、通信和安全约束。这样的混合式架构能够降低纯强化学习试错带来的工程风险，也更容易解释决策依据。

最终，多无人机协同需要以任务完成结果评价，而不是以保持队形或局部轨迹一致作为唯一指标。对于工业巡检，真正重要的指标包括关键设备覆盖率、异常发现率、异常复查时延、任务完成时间、单位任务能耗以及通信开销。只有建立这些面向巡检任务的指标，才能判断一种多机协同算法是否真正提高了工业巡检效率和可靠性。

\section{总结与展望}

\subsection{现有研究存在的问题}

综合国内外研究可以看出，复杂工业场景无人机自主巡检已经形成较为完整的技术体系。环境感知方面，视觉、激光和惯性融合已经能够在部分GNSS拒止和受限空间中实现稳定定位；异常认知方面，监督学习、无监督异常检测和基础模型不断提高对已知与未知缺陷的识别能力；自主决策方面，覆盖路径规划、动态重规划以及多无人机协同逐步成熟，部分研究已经进入真实飞行验证阶段。但是，从“能够完成一次巡检”走向“能够在真实工业场景中长期、安全、稳定地自主运行”，仍存在明显差距。综合前述研究，可将当前主要问题概括为以下五个方面。

（1）复杂退化环境下的自主感知鲁棒性仍然不足。现有SLAM和多传感器融合方法主要围绕精度和实时性进行优化，而工业现场中的强反光、弱纹理、粉尘烟雾、动态遮挡、GNSS多径以及电磁干扰往往会同时影响多个传感器。当前不少方法在单一退化因素下已经具有一定适应能力，但对于多因素叠加造成的观测异常和传感器失效，仍缺少成熟的在线识别与自适应融合机制
\cite{cadena2016past,lvyang2019sense,gupta2022slam,petrlik2020robust}。因此，下一阶段感知研究不能仅强调更低的定位误差，而应进一步研究“感知可靠度”的估计，使系统能够知道自己什么时候可信、什么时候需要切换感知模式。

（2）工业异常认知仍存在跨场景泛化和未知异常识别能力不足的问题。监督式检测在特定设备和缺陷类别上可以达到很高精度，但高度依赖标注数据，且模型通常与设备型号和成像条件绑定；无监督异常检测减少了标注需求，却容易把正常工况变化、光照变化和视角变化误判为异常
\cite{taox2021survey,roth2022towards,batzner2024efficientad}. 基础模型和视觉语言模型为零样本、少样本和开放词汇识别提供了新途径，但目前多数验证仍集中在公开数据集，真实工业场景下的领域适配、异常解释和不确定性量化还不充分
\cite{jeong2023winclip,kirillov2023segment,liu2024grounding,oquab2024dinov2}。因此，未来需要研究能够同时利用视觉、红外、三维几何和设备知识的异常认知体系。

（3）自主巡检规划与设备状态之间的耦合仍然不够。现有路径规划研究已经能够很好地解决从起点到终点、覆盖区域和动态避障等问题，但传统目标函数主要围绕路径长度、时间或避障安全建立。实际工业巡检的任务收益却由“有没有发现重要异常、异常是否需要立即复查、观测质量是否足够”决定
\cite{choset2001coverage,galceran2013survey,bircher2016receding}. 如果规划系统不理解设备状态，那么即使路径最优，也可能错过真正需要关注的对象。因此，未来需要把异常等级、设备风险和观测不确定性直接纳入任务规划和重规划过程。

（4）多无人机协同的真实工业应用仍不充分。现有研究已经从集中式任务分配发展到分布式协同与多智能体强化学习，并在野外和实验环境中展示了很强的自主能力
\cite{choi2009consensus,gronauer2022marl,zhou2022swarm,zhou2023racer}. 但真实工业环境中还存在通信带宽受限、链路中断、异构载荷、动态任务和安全隔离等问题。尤其在异常事件驱动的任务切换方面，如何保证多无人机在有限通信条件下迅速形成新的任务分工，仍缺乏统一而可靠的解决方案。

（5）现有研究尚未真正形成贯通感知、认知和决策的闭环系统。许多工作分别追求SLAM定位精度、异常检测准确率或路径规划效率，各模块之间的信息交互较弱。例如异常检测结果通常只保存为一条报警记录，没有直接改变后续飞行路径；路径规划也很少利用模型的不确定性决定下一次观测位置。实际上，自主巡检的核心应是一个持续闭环：“感知环境--理解设备--判断风险--决定动作--执行任务--再次感知”。只有把模块评价从单点指标提高到任务闭环指标，才能真正衡量无人机系统的自主程度。

此外，还有一个贯穿上述问题的工程瓶颈，即算法从实验室向工业现场迁移的鸿沟。实验室研究通常使用公开数据集、理想照明和稳定通信，能够较容易获得清晰的性能提升；真实工业现场则存在大量不可控因素，而且安全测试成本高、数据获取难、异常事件不能人为重复。工业自主巡检要达到长期运行要求，还需要解决模型更新、数据闭环、设备台账关联、软件故障恢复以及系统安全验证等问题。因此，未来研究不仅需要更先进的算法，还需要更加完整的系统工程方法和长期实测机制。

\subsection{未来发展趋势}

（1）面向退化环境的多模态鲁棒感知将成为自主巡检的基础能力。未来感知系统应由“固定传感器组合”向“可感知退化状态、可动态调整融合权重”的自适应体系发展。视觉、激光、红外、IMU和GNSS可以在不同运行阶段承担不同作用，当视觉受到烟雾或低照度影响时，系统可以提高激光和惯性信息权重；当点云稀疏或局部反射异常时，则可以利用视觉和历史地图进行补充。更进一步，可以建立传感器健康状态和观测可信度模型，使系统在某个模态失效时能够自动降级运行而不是整体失效。复杂工业场景中的鲁棒性，最终应体现为“传感器坏一个、环境变一种，系统仍然可以安全工作”。

（2）基础模型将推动工业异常认知由封闭类别向开放世界演进。传统检测器通常需要预先定义类别和缺陷标签，而基础模型可以利用自然语言先验和大规模预训练知识，为新设备和新缺陷提供迁移能力
\cite{radford2021learning,jeong2023winclip,oquab2024dinov2}. 未来可以将设备名称、部件名称、异常描述以及维修知识作为语义先验，与视觉和三维数据联合使用，使模型能够从“检测某个已知类别”进一步发展为“解释一个异常现象”。但在工业安全关键领域，基础模型的输出仍需结合确定性规则、物理约束和多传感器证据，建立可验证的异常判断机制，而不能完全依赖黑盒生成结果。

（3）感知驱动的主动巡检将成为区别于传统自动航线的重要方向。未来无人机不应只执行预先规划的任务，而应根据当前设备状态动态决定观察位置、观察次数和观测模态。当第一次巡检发现疑似腐蚀后，系统应能够自动选择更合适的距离和角度进行复检；当红外发现热异常时，可以触发可见光近距离观测以寻找对应结构缺陷；当多个设备同时异常时，可以根据风险等级和剩余能量重新排序任务。这类“发现异常--主动复查--确认状态”的闭环能够显著提高有限飞行时间内的信息利用率，并使无人机真正从“数据采集器”变为“主动巡检执行者”。

（4）数字孪生将成为巡检数据从一次性记录走向长期状态管理的重要载体。无人机能够提供高频率空间观测，但如果缺乏统一设备对象和历史状态管理，数据往往仍然是孤立图像和视频。将无人机图像、点云、温度和异常结果映射到设备数字孪生模型，可以在统一坐标中记录设备结构、历史缺陷和状态变化，为缺陷演化分析和预测性维护提供条件
\cite{tao2019digital}. 更进一步，数字孪生不仅接收巡检结果，也可以反过来生成巡检需求，例如根据设备健康度决定某个区域是否需要增加巡检频率，从而形成“状态驱动巡检”的闭环。

（5）多无人机协同将由静态分工走向风险驱动的动态协同。传统多机任务分配主要根据距离、任务量和剩余能量进行，而未来应进一步引入设备风险等级、异常置信度、任务时效以及载荷能力等因素。不同无人机可以根据自身能力承担不同子任务，形成发现、复检、测温和三维重建等角色分工。当通信受限时，系统还应能够根据局部信息完成局部决策，并在通信恢复后进行全局一致性修正。多智能体强化学习、分布式优化和任务逻辑约束有望在这一过程中发挥作用，但需要与工程安全规则融合，以保证策略可解释和可验证
\cite{yu2022surprising,zhou2023racer}。

（6）端边云协同和算力感知将成为自主巡检系统规模化部署的重要条件。无人机机载计算资源有限，不适合无限制运行大型基础模型，但将所有数据全部上传到云端又会受到通信带宽和时延约束。因此，未来可以将任务分为机载快速筛查、边缘服务器精细分析和云端历史建模三个层次。机载侧完成定位、障碍感知和轻量异常筛查；厂区边缘节点负责高分辨率视觉、三维重建和多模态联合推理；云端则完成跨时间尺度的数据管理、模型更新和全局任务优化
\cite{zhou2018computation,mozaffari2019tutorial}. 这种计算架构能够在保证实时性的同时降低无人机本体的能源负担。

总体而言，复杂工业场景无人机自主巡检的发展方向并不是简单追求某一种算法在某个数据集上的最高指标，而是形成能够长期运行的系统级自主能力。其技术闭环可以概括为“多模态感知--自主定位--设备认知--风险评估--主动规划--任务执行--状态反馈”，并逐步与多无人机协同、数字孪生和端边云计算融合。未来的评价标准也应从单模块精度逐步转向任务完成率、异常发现时效、定位可靠性、误报漏报代价、能源利用率以及长期运行稳定性。对于真正面向工业现场的研究而言，算法创新只有在能够经受真实环境、真实任务和长期运行的检验时，才能最终转化为具有工程价值的自主巡检能力。

